\documentclass[10pt,a4paper]{amsart}
\usepackage[margin=19mm]{geometry}
\usepackage{amsmath,amssymb,mathtools}
\usepackage[scr=rsfs,cal=euler]{mathalfa}
\usepackage{xfrac}
\usepackage[unicode,hidelinks,bookmarks=false]{hyperref}
\newcommand{\ie}{i.e.\ }
\newcommand{\cf}{cf.\ }
\newcommand{\resp}{respectively\ }
\newcommand{\Subsection}{\subsection}
\providecommand{\nobreakdash}{\nobreak\hbox{-}}
\setlength{\emergencystretch}{2em}
\multlinegap0pt
\usepackage{amssymb}
\usepackage{tikz}
\usepackage{xcolor}
\usepackage{mathtools}
\newtheorem{theorem}{Theorem}
\newtheorem{claim}{Claim}
\usepackage{cleveref}
\crefname{theorem}{Theorem}{Theorems}
\crefname{claim}{Claim}{Claims}
\DeclareMathOperator{\ur}{srr}
\title{The Log-Rank Conjecture: New Equivalent Formulations}
\author{Lianna Hambardzumyan, Shachar Lovett, Morgan Shirley}
\date{Source: arXiv:2510.02583v3 (CC BY 4.0)}
\begin{document}
\maketitle

\noindent\emph{Source and licence.} The Log-Rank Conjecture: New Equivalent Formulations. Version arXiv:2510.02583v3 (CC BY 4.0), distributed by arXiv under the Creative Commons Attribution 4.0 International licence, http://creativecommons.org/licenses/by/4.0/, as recorded in the arXiv licence metadata for this exact version and shown on the abstract page https://arxiv.org/abs/2510.02583v3. Full licence notice: \texttt{licenses/arxiv-2510.02583-CC-BY-4.0.txt}.\par\smallskip

\noindent\textit{Editorial scope.} Counted unit: the article's theorem thm:main (source0.tex lines 117--120). The article's complete original proof of it is delivered with it (source0.tex lines 363--367, one \texttt{proof} environment of 5 lines, which the article places away from the statement). No mathematics is written, repaired or re-derived: the statement and the proof are the article's own lines, in the article's own notation, and no retained line is substituted or re-flowed.  Retained uncounted as the source-local context the proof needs: lines 334--336; lines 348--350; lines 356--358.  Nothing was omitted from any retained span, so the package carries no omission record.  No retained line contains a citation, so the wrapper carries no bibliography and no bibliography entry.  No retained line contains a figure environment, an image-inclusion command or drawing code, so no image is delivered and the package declares no asset dependency.  Editorial formatting only, in addition: an AMS \texttt{amsart} preamble that restates the article's own declarations and macros used by the retained lines, namely the environments \texttt{theorem}, \texttt{claim} on separate counters, together with the standard packages those lines need.  The retained environments print in this document as Theorem 1, Claim 1 in order of appearance, whereas the article prints the same items with the numbering of its own full document.  The complete native source of the article as distributed in the arXiv e-print for this exact version is bundled unchanged as \texttt{sources/186292/source0.tex}.
\begin{theorem} \label{thm:tensors}
	Let $T: [n_1] \times \ldots \times [n_\ell] \to \{0, 1\}$ be an order-$\ell$ Boolean tensor with $\ell \geq 2$. If the tensor rank of $T$ is $r$, then $T$ can be written as a $\pm 1$-linear-combination of at most $(c r \log r)^{(\ell - 1)}$ primitive tensors, where $c$ is an absolute constant.
\end{theorem}

\begin{claim} \label{cl:IS-tensors}
	Let $S$ be an independent set of $\lambda$-slices of $T$. Then $|S| \leq O(r \log r)$, where $r$ is the tensor rank of $T$.
\end{claim}

\begin{claim}\label{cl:tensor_linear_combin}
	Let $S$ be a maximal independent set of $\lambda$-slices of $T$.  Then every $\lambda$-slice of $T$ can be expressed as a  $\pm 1$-linear combination of  $\lambda$-slices in $S$.
\end{claim}

\begin{theorem}[Main Theorem]\label{thm:main}
	Every Boolean matrix $M$ of rank $r$ can be written as a $\pm 1$-linear-combination of at most $O(r \log r)$ primitive matrices, that is,
	\[\ur(M) \leq O(r \log r).\]
\end{theorem}

\begin{proof}[Proof of \cref{thm:tensors}]
	The proof is by induction. The base case is $\ell=2$, which is proven by \cref{thm:main}. For $\ell>2$, arbitrarily choose a coordinate  $\lambda$ and use \cref{cl:tensor_linear_combin} to express $T$ as a $\pm 1$-linear combination of a maximal independent set of $\lambda$-slices $S$. Each $\lambda$-slice of $S$ has tensor rank at most $r$, and so by the inductive hypothesis, each $\lambda$-slice in $S$ can be expressed as a $\pm 1$-linear combination of $(cr \log r)^{(\ell - 2)}$ order-$(\ell-1)$ primitive tensors.
	
	In a similar fashion to the proof of \cref{thm:main}, we therefore can write $T$ as a $\pm 1$-sum of ${|S| \cdot (cr \log r)^{(\ell - 2)}}$ order-$\ell$ primitive tensors. Conclude by substituting $|S| = O(r \log r)$ using \cref{cl:IS-tensors}.
\end{proof}

\end{document}
